\section*{Capítulo \ref{ch:b2:coordination-complexes} --- Metales de transición y complejos de coordinación}

\begin{solution}{exo:b2:coordination-complexes:1}
\ce{Ti^3+} $d^1$, \ce{V^3+} $d^2$, \ce{Cr^3+} $d^3$, \ce{Mn^2+} $d^5$, \ce{Fe^3+} $d^5$, \ce{Co^2+}
$d^7$, \ce{Ni^2+} $d^8$, \ce{Zn^2+} $d^{10}$.
\end{solution}

\begin{solution}{exo:b2:coordination-complexes:2}
Agua, 1; en, 2; etanodioato, 2; edta, 6; tiocianato, 1; el tiocianato es ambidentado
(por S o por N).
\end{solution}

\begin{solution}{exo:b2:coordination-complexes:3}
Sulfato de tetraamminocobre(II); hexacianidoferrato(III) de potasio;
cloruro de tetraamminodicloridocobalto(III); tetracarbonilníquel(0).
\end{solution}

\begin{solution}{exo:b2:coordination-complexes:4}
Lineal; tetraédrica; plano-cuadrada ($d^8$, tercera serie); octaédrica.
\end{solution}

\begin{solution}{exo:b2:coordination-complexes:5}
Dos: cis (los dos cloruros adyacentes) y trans (opuestos). En un tetraedro
solo habría uno.
\end{solution}

\begin{solution}{exo:b2:coordination-complexes:6}
\ce{Cr(CO)6} 18; \ce{Fe(CO)5} 18; \ce{Ni(CO)4} 18; \ce{V(CO)6}: $5 + 12 = 17$;
\ce{Mn2(CO)10}, 18 por manganeso. \ce{V(CO)6} incumple la regla: es \omterm{def:b2:diatomic-mos:magnetism}{paramagnético} y
se reduce con facilidad a \ce{[V(CO)6]-}, que tiene 18.
\end{solution}

\begin{solution}{exo:b2:coordination-complexes:7}
Ferroceno: Fe(II) $d^6$ + 12 = 18. Cobaltoceno: Co(II) $d^7$ + 12 = 19. El cobaltoceno pierde
con facilidad su electrón adicional y da el ion cobaltocenio \ce{[Co(C5H5)2]+}, con 18.
\end{solution}

\begin{solution}{exo:b2:coordination-complexes:8}
Un edta sustituye a seis moléculas de agua: \ce{[Ni(H2O)6]^2+ + edta^4- -> [Ni(edta)]^2- +
6H2O}, siete partículas a partir de dos. La \omterm{def:b2:reaction-free-energy:reaction-entropy}{entropía de reacción} es fuertemente positiva, y
$K$, muy grande. Seis ligandos independientes sustituyen a seis moléculas de agua sin ganancia en el
número de partículas.
\end{solution}

\begin{solution}{exo:b2:coordination-complexes:9}
Nitrito-$\kappa N$ (\ce{Co-NO2}, el isómero nitro, amarillo) y nitrito-$\kappa O$
(\ce{Co-O-N=O}, el isómero nitrito, rojo).
\end{solution}

\begin{solution}{exo:b2:coordination-complexes:10}
Tres: el isómero trans (aquiral: planos de simetría) y los dos enantiómeros del isómero
cis (los dos quelatos y los dos cloruros se enrollan en un sentido o en el otro).
\end{solution}

\begin{solution}{exo:b2:coordination-complexes:11}
Hexágono plano: B en las posiciones 1,2 (orto), 1,3 (meta), 1,4 (para): tres isómeros.
Prisma trigonal: B a lo largo de una arista de un triángulo, a lo largo de una arista vertical, o en diagonal en una
cara rectangular: tres. El octaedro da dos. No encontrar más de dos, pese a
muchos intentos y a varios compuestos, descarta las otras dos geometrías, siempre que no se haya
pasado por alto un tercer isómero.
\end{solution}

\begin{solution}{exo:b2:coordination-complexes:12}
\ce{[RhCl(PPh3)3]}: Rh(I) $d^8$ + 4 × 2 = 16. \ce{[IrH(CO)(PPh3)3]}: Ir(I) $d^8$ + 5 × 2 =
18. El complejo de rodio, con 16 electrones, puede admitir dos más (adición oxidante,
\cref{ch:b2:catalytic-cycles}).
\end{solution}

\begin{solution}{pb:b2:coordination-complexes:1}
\textbf{1.} Los cloruros fuera de la \omterm{def:b2:coordination-complexes:coordination-sphere}{esfera de coordinación}, libres en disolución.
\textbf{2.} A: 0; B: 1; C y D: 2.
\textbf{3.} A: $[\ce{Co(NH3)6}]^{3+}$ + 3 \ce{Cl-}: 4 iones; B: 3; C, D: 2. Coinciden.
\textbf{4.} +3.
\textbf{5.} $d^6$.
\textbf{6.} Seis en cada uno: 6 \ce{NH3}; 5 \ce{NH3} + 1 Cl; 4 \ce{NH3} + 2 Cl.
\textbf{7.} \ce{[Co(NH3)6]Cl3}; \ce{[CoCl(NH3)5]Cl2}; \ce{[CoCl2(NH3)4]Cl} (dos veces).
\textbf{8.} Cloruro de hexaamminocobalto(III); cloruro de pentaamminocloridocobalto(III).
\textbf{9.} Tetraamminodicloridocobalto(III).
\textbf{10.} Isómeros geométricos (cis--trans).
\textbf{11.} \ce{[CoCl3(NH3)3]}, un complejo neutro: no da iones ni precipita cloruro de
inmediato.
\textbf{12.} Dos: fac y mer.
\textbf{13.} Cis: los dos Cl a $90^\circ$; trans: a $180^\circ$.
\textbf{14.} Tres (posiciones orto, meta y para en el hexágono).
\textbf{15.} Tres.
\textbf{16.} Apoya el octaedro, la única de las tres geometrías que predice dos.
\textbf{17.} No: la ausencia de un tercer isómero es una prueba negativa. Un tercer isómero
habría descartado el octaedro.
\textbf{18.} El que forma el quelato: un etanodioato bidentado solo puede abarcar dos
posiciones cis.
\textbf{19.} Tres quelatos en alrededor del cobalto, cada uno abarcando dos posiciones cis.
\textbf{20.} Ninguno de los dos: no hay plano de simetría ni centro de inversión.
\textbf{21.} Dos enantiómeros, $\Delta$ y $\Lambda$.
\textbf{22.} Que los seis ligandos están dispuestos en tres dimensiones, como en un
octaedro; una disposición plana daría un ion aquiral.
\textbf{23.} Cis: quiral (dos enantiómeros). Trans: aquiral.
\textbf{24.} Mostró que la actividad óptica es una propiedad de la geometría molecular en general,
no del carbono, y confirmó el octaedro.
\textbf{25.} El octaedro predice $\boldsymbol{2}$ isómeros de \ce{[CoCl2(NH3)4]+}; el
hexágono plano y el prisma trigonal predicen $\boldsymbol{3}$.
\end{solution}
